\addtocontents{toc}{\protect\Needspace{3\baselineskip}}
\originalchapter{场的维数与引力的长度尺度}\label{ch:basic-fields}
对应 8.251 第 3--4 讲. 电磁场为如何把普通物理方程写成相对论张量提供例子, 高维高斯定律则把额外维度从图形识别变成可计算的作用力. 本章只使用弱场与静态引力; 广义相对论的完整曲率推导以前置的 GR 卷为准.

\originalsection{势、规范自由度与场强}
采用原课的有理化单位, 使 $\nabla\cdot\boldsymbol E=\rho$ 且 $\boldsymbol E,\boldsymbol B$ 同量纲. 在三维空间
\[
 \boldsymbol B=\nabla\times\boldsymbol A,\qquad
 \boldsymbol E=-\nabla\Phi-\frac1c\partial_t\boldsymbol A.
\]
规范变换 $\boldsymbol A\mapsto\boldsymbol A+\nabla\epsilon$,
$\Phi\mapsto\Phi-c^{-1}\partial_t\epsilon$ 不改变这两个场. 把势组合为 $A^\mu=(\Phi,\boldsymbol A)$, 则 $A_0=-\Phi$, 并定义
\[
 F_{\mu\nu}=\partial_\mu A_\nu-\partial_\nu A_\mu,
 \quad F_{0i}=-E_i,\quad F_{ij}=\epsilon_{ijk}B_k.
\]
由于偏导交换, $A_\mu\mapsto A_\mu+\partial_\mu\epsilon$ 使 $\delta F_{\mu\nu}=0$. 其反对称性提供 $D(D-1)/2$ 个独立分量, 而非对称张量的 $D(D+1)/2$ 个.

在单连通、拓扑平凡的适当区域内, 同一场强的两个势局部相差一个梯度. 全局情形可能存在不能这样消除的环路积分, 所以“场强相等”与“全局规范等价”应分别检查. 这一区分在圆上的 Wilson 环尤其重要.

\originalsection{麦克斯韦方程怎样装进张量式}
令 $J^\mu=(c\rho,\boldsymbol J)$. 方程为
\[
 \partial_\nu F^{\mu\nu}=\frac{J^\mu}{c},\qquad
 \partial_\lambda F_{\mu\nu}+\partial_\mu F_{\nu\lambda}
 +\partial_\nu F_{\lambda\mu}=0.
\]
第一式的 $\mu=0$ 分量给 $\partial_iF^{0i}=\nabla\cdot\boldsymbol E=\rho$. 空间分量则给 $\nabla\times\boldsymbol B-c^{-1}\partial_t\boldsymbol E=\boldsymbol J/c$. 第二式用 $F=\dd A$ 展开后各二阶偏导两两抵消; 在三维空间分别还原 $\nabla\cdot\boldsymbol B=0$ 与法拉第感应定律. 对第一式再取 $\partial_\mu$ 得 $\partial_\mu J^\mu=0$, 因对称的双偏导与反对称的 $F$ 缩并为零.

在 $d$ 个空间维度, 电场仍是 $d$ 分量的向量, 磁场则一般是空间反对称二阶张量, 有 $d(d-1)/2$ 个分量. 只有 $d=3$ 时它可由 $\epsilon_{ijk}$ 对偶成另一个三分量向量. 对 $d=2$, 唯一的 $F_{12}$ 可表示为一个赝标量. 把三维叉乘照抄进四维空间会漏掉磁场分量; 张量方程不受这个限制.

\originalsection{高斯定律与单位球面的面积}
记单位 $(d-1)$ 维球面面积为
\[
 \Omega_{d-1}=\frac{2\pi^{d/2}}{\Gamma(d/2)},\qquad
 \Gamma(a)=\int_0^\infty t^{a-1}\ee^{-t}\dd t.
\]
用 $d$ 维高斯积分可直接求它: 笛卡尔坐标给 $\int\ee^{-|x|^2}\dd^dx=\pi^{d/2}$, 球坐标给 $\Omega_{d-1}\int_0^\infty r^{d-1}\ee^{-r^2}\dd r=\Omega_{d-1}\Gamma(d/2)/2$. 两者相等即上式, 且 $\Omega_1=2\pi,\Omega_2=4\pi,\Omega_3=2\pi^2$.

对原点电荷 $q$, 球对称与散度定理给
\[
 E(r)\Omega_{d-1}r^{d-1}=q,
 \qquad E(r)=\frac{q}{\Omega_{d-1}r^{d-1}}.
\]
因此电场幂次来源于通量穿过球面的面积增长, 不是预先选定的反平方律. $d>2$ 时可取无穷远势为零,
\[
 \Phi(r)=\frac{q}{(d-2)\Omega_{d-1}r^{d-2}};
\]
$d=2$ 时则 $\Phi=-q\ln(r/r_0)/(2\pi)$, 无法同时要求对数势在无穷远为零. 此处 $S^{d-1}$ 表示球面本身, 维数为 $d-1$; 原笔记把个别 $S^0,S^1$ 的维数混写, 本稿按此统一.

\originalsection{高维牛顿引力及轨道稳定性}
按原课约定在 $D=d+1$ 维写 $\nabla^2V_D=4\pi G_D\rho_D$, 引力加速度 $\boldsymbol g=-\nabla V_D$. 对点质量 $M$, $d>2$ 时
\[
 V_D=-\frac{4\pi G_DM}{(d-2)\Omega_{d-1}r^{d-2}},\qquad
 g_r=-\frac{4\pi G_DM}{\Omega_{d-1}r^{d-1}}.
\]
在 $d=3$ 恢复 $V=-GM/r$. 不同文献常把 $\Omega_{d-1}$ 吸收入高维 $G_D$, 比较额外维度关系前必须统一其泊松方程定义.

为检查“高维轨道不稳定”的条件, 设中心势 $V=-A/r^q$, $A>0,q=d-2>0$, 角动量大小为 $\ell$. 单位粒子势能 $mV$ 与离心项形成
\[
 U_{\rm eff}=\frac{\ell^2}{2mr^2}-\frac{mA}{r^q}.
\]
在圆轨道处 $\ell^2/m=mAqr^{2-q}$, 代入二阶导数得 $U_{\rm eff}''=mAq(2-q)/r^{q+2}$. 所以 $d=3$ 的圆轨道径向稳定, $d=4$ 在纯幂势下为边界情形, $d>4$ 不稳定. 这不是对所有高维物理模型的生命存在性证明; 紧致化后远距离势也会重新变成三维形式.

\originalsection{普朗克单位与弱场规范变换}
由 $G,c,\hbar$ 构造长度、质量、时间,
\[
 \ell_P=\sqrt{\frac{G\hbar}{c^3}},\qquad
 m_P=\sqrt{\frac{\hbar c}{G}},\qquad t_P=\ell_P/c.
\]
采用近似数值 $G=6.6743\times10^{-11}\,\mathrm{m^3kg^{-1}s^{-2}}$,
$c=2.99792458\times10^8\,\mathrm{m\,s^{-1}}$,
$\hbar=1.054571817\times10^{-34}\,\mathrm{J\,s}$,
有 $\ell_P\simeq1.6163\times10^{-35}\,\mathrm m$,
$m_P\simeq2.1764\times10^{-8}\,\mathrm{kg}$,
$t_P\simeq5.3912\times10^{-44}\,\mathrm s$.
这些组合给量子引力的自然尺度, 并不单独证明时空在该尺度离散化.

弱引力写 $g_{\mu\nu}=\eta_{\mu\nu}+h_{\mu\nu}$. 对被动坐标变换 $x^{\prime\mu}=x^\mu+\xi^\mu(x)$, 线性阶
\[
 h'_{\mu\nu}=h_{\mu\nu}-\partial_\mu\xi_\nu-\partial_\nu\xi_\mu.
\]
若改用主动变换或 $x'=x-\xi$, 两项的符号相反. 这与电磁势的规范变换类似, 但引力有两个张量指标. 准确的非线性几何在 GR 卷介绍, 本卷不再重推爱因斯坦方程.

用 $m_P/\ell_P^3$ 估算的密度约为 $5.2\times10^{96}\,\mathrm{kg\,m^{-3}}$, 与宇宙学暗能量密度有巨大差距. 原课给出的 2007 年数量级差约 $10^{122}$, 是关于自然尺度的对比, 不是严格推导真空能必等于该值. 更精细的场论估计还取决于截止、对称性和重整化条件.

\originalsection{紧致化如何改变有效引力常数}
考虑一个周长 $L=2\pi R$ 的紧致方向 $y$. 对在 $y$ 上均匀的零模式源,
$\rho_5(\boldsymbol x,y)=\rho_4(\boldsymbol x)/L$. 此时 $\partial_yV=0$, 五维泊松方程变成
\[
 \nabla_3^2V=4\pi\frac{G_5}{L}\rho_4,
 \qquad G_4=\frac{G_5}{L}.
\]
对局域源, 远距离 $r\gg R$ 仍由零模式主导, 非零 Kaluza--Klein 模产生指数衰减修正. 若紧致空间体积为 $V_n$、几何与度规取最简单直积, 则 $G_4=G_{4+n}/V_n$. 翘曲、非平凡零模归一化和额外轻场需要重新积分, 不能无条件使用此体积关系.

\originalsection{大额外维度模型的数量级}
沿原课约定定义 $G_D=(\ell_{P,D})^{D-2}c^3/\hbar$, 令 $n=D-4$ 个紧致方向周长均为 $L$. 有
\[
 L^n=\frac{(\ell_{P,D})^{n+2}}{\ell_P^2}.
\]
用原课的数量级 $\ell_{P,D}=10^{-18}\,\mathrm{cm}$,
$\ell_P=10^{-33}\,\mathrm{cm}$, $n=1$ 给 $L=10^{12}\,\mathrm{cm}$, $n=2$ 给 $L=10^{-3}\,\mathrm{cm}=10\,\mu\mathrm m$. 这里 $L$ 是周长, 半径还需除 $2\pi$. 原笔记写 $1\,\mu\mathrm m=10^{-6}\,\mathrm{cm}$ 是单位笔误, 应为 $10^{-4}\,\mathrm{cm}$.

若标准模型场也在大额外维度自由传播, 会产生容易观察到的低质量紧致模. 因此这类模型常假设通常物质场局限在 D 膜, 而引力进入体空间. 这是模型假设, 不能写成已实验证实“所有电子都是开弦”的结论. 原课的加速器与短程引力实验叙述属于 2007 年背景; 本稿不把该历史数值当作当前排除界限.
